\documentclass[11pt]{article}
\usepackage[margin=1in]{geometry}
\usepackage{hyperref}
\usepackage{enumitem}
\usepackage{titlesec}
\usepackage{graphicx}
\usepackage{amsmath}
\usepackage{booktabs}

\title{\textbf{HarborGuard AI / Datakly}\\
Autonomous Maritime Disruption Optimization Platform}
\author{Potential HTF Idea}
\date{}

\begin{document}

\maketitle

\section*{1. Executive Summary}

HarborGuard AI is an AI-native maritime disruption decision-support platform built for U.S.-based mid-market electronics manufacturers ($100M$--$300M$ revenue).

Most mid-market manufacturers lose \textbf{10--15\% of revenue annually} due to preventable supply chain disruptions. Not because they lack data, but because they lack rapid, structured, financially-optimized decision intelligence.

HarborGuard AI acts as an operational control tower powered by a risk-aware optimization engine. It continuously monitors maritime disruption signals (e.g., Port of Los Angeles), models exposure across inbound SKUs, simulates mitigation trade-offs, and recommends financially optimal strategies aligned with a company’s risk DNA.

This is a SaaS decision-support system (not full autonomy) with human-in-the-loop safeguards.

\section*{2. Target Customer (ICP)}

\textbf{Ideal Customer Profile (ICP):}

\begin{itemize}
    \item U.S.-based electronics manufacturer
    \item Revenue: $120M$--$250M$
    \item $\geq$ 40\% of components imported from Asia
    \item $\geq$ 35\% inbound concentration through Port of Los Angeles
    \item SLA commitments $\geq$ 95\%
    \item No dedicated supply chain risk intelligence team
\end{itemize}

\textbf{Primary Buyer:} VP of Supply Chain  
\textbf{Economic Sign-off:} COO or CFO  

\section*{3. Problem Statement}

When a maritime disruption occurs:

\begin{itemize}
    \item Teams scramble across ERP dashboards and spreadsheets
    \item Decisions are delayed
    \item Trade-offs between cost and service are unclear
    \item Revenue-at-risk accumulates quickly
\end{itemize}

Existing systems (SAP, Oracle, visibility tools) provide tracking and dashboards. They do not provide proactive, risk-aware trade-off optimization.

\section*{4. Product Overview}

\textbf{Product Type:} SaaS Platform (Decision-Support Only)

HarborGuard AI provides:

\begin{enumerate}
    \item Continuous maritime risk monitoring
    \item Exposure mapping to inbound SKUs
    \item Risk-DNA-aware trade-off optimization
    \item Financial impact simulation
    \item Autonomous mitigation drafting
    \item Explainable reasoning
\end{enumerate}

\section*{5. Core Innovation: Risk DNA Optimization Engine}

Each manufacturer has a structured Risk DNA profile:

\begin{itemize}
    \item Risk tolerance coefficient
    \item SLA weight
    \item Working capital constraint
    \item Customer concentration sensitivity
    \item SKU-level margin variability
\end{itemize}

The system solves:

\[
\max (Revenue\ Preserved - Mitigation\ Cost - SLA\ Penalty)
\]

Subject to:

\begin{itemize}
    \item Inventory depletion constraints
    \item Lead-time uncertainty
    \item Working capital cap
    \item Risk tolerance boundary
\end{itemize}

The engine evaluates strategies:

\begin{itemize}
    \item No action
    \item Full air freight
    \item Reroute to alternate port
    \item Split strategy (reroute + air + buffer stock)
\end{itemize}

\section*{6. System Architecture}

\subsection*{6.1 Perception Layer}
\begin{itemize}
    \item News ingestion (LLM classification)
    \item Port congestion monitoring
    \item Disruption severity scoring
\end{itemize}

\subsection*{6.2 Exposure Mapping}
\begin{itemize}
    \item Map inbound shipments to port
    \item SKU criticality analysis
    \item Inventory runway calculation
\end{itemize}

\subsection*{6.3 Optimization Engine}
\begin{itemize}
    \item Multi-strategy simulation
    \item Revenue-at-risk estimation
    \item Cost-impact modeling
\end{itemize}

\subsection*{6.4 Action Layer}
\begin{itemize}
    \item Supplier rerouting email drafts
    \item Logistics rebooking requests
    \item Executive escalation memo
\end{itemize}

\subsection*{6.5 Memory Layer}
\begin{itemize}
    \item Logs past disruptions
    \item Tracks mitigation success
    \item Adjusts risk weights
\end{itemize}

\section*{7. Demo Scenario}

Simulated event:

Labor strike at Port of Los Angeles  
10--14 day disruption  

Output:

\begin{itemize}
    \item 47 inbound containers affected
    \item 18 high-priority SKUs at risk
    \item Revenue-at-risk (45 days): \$38.4M
\end{itemize}

Recommended Strategy:

\begin{itemize}
    \item Reroute 60\%
    \item Air freight 20\%
    \item Buffer 20\%
\end{itemize}

Financial Impact:

\begin{itemize}
    \item Mitigation cost: \$2.8M
    \item Revenue preserved: \$34.1M
    \item Gross margin protected: \$9.5M
\end{itemize}

\section*{8. Competitive Positioning}

Competitors:
\begin{itemize}
    \item SAP IBP
    \item Oracle SCM Cloud
    \item project44
    \item FourKites
\end{itemize}

Differentiation:

\begin{itemize}
    \item Proactive optimization vs passive visibility
    \item Risk-DNA-aware decisions
    \item Financially grounded trade-off modeling
\end{itemize}

\section*{9. Pricing Strategy}

Base: \$90K/year  
Growth: \$140K/year  

ROI Example:

For \$200M manufacturer:

\begin{itemize}
    \item 8\% disruption volatility = \$16M exposure
    \item 25\% impact reduction = \$4M preserved
    \item 28\% margin = \$1.12M gross profit saved
\end{itemize}

9x ROI on subscription.

\section*{10. Responsible AI Safeguards}

\begin{itemize}
    \item Human approval required before execution
    \item Confidence intervals shown
    \item Explainability panel
    \item Clear override capability
\end{itemize}

\section*{11. 1-Week Build Plan}

\subsection*{Engineering Needs}

Perception + Disruption Scoring |
Exposure Mapping + Inventory Modeling |
Optimization Engine + LLM Action Layer  

Deliverables:
\begin{itemize}
    \item Working simulation
    \item Risk scoring
    \item Strategy comparison output
    \item Auto-generated mitigation drafts
\end{itemize}

\subsection*{Business Side}

\begin{itemize}
    \item Risk DNA modeling framework
    \item Financial impact logic
    \item ROI modeling
    \item Pitch narrative
    \item ------------------------------
    \item Go-to-market strategy
    \item Pricing structure
    \item Competitive analysis
    \item VC roadmap framing
\end{itemize}

\section*{12. Venture-Scale Vision}

Phase 1: Maritime Disruption Optimization  \\
Phase 2: Tariff + Geopolitical Risk Modeling  \\
Phase 3: Supplier Insolvency Intelligence  \\
Phase 4: Full Disruption Intelligence Platform  \\

\end{document}